\begin{enumerate}[label=\color{red}\arabic*), leftmargin=*]
	\item \lb{Determina cuáles de los siguientes subconjuntos de $\R^3$ son subespacios vectoriales:}
	\begin{enumerate}[label=\color{red}\alph*)]
		\item $\db{\{(x,y,z)\in\R^3:x=1\}}$ 
		
		No porque el vector $(1,0,0), (1,1,0)$ pertenece al subespacio pero $(1,0,0)+(1,1,0)=(2,1,0)$ no pertenece.
		\item $\db{\{(x,y,z)\in\R^3:xyz=0\}}$
		
		No porque al comprobar los la suma de vectores $(1,0,0)+(0,0,1)=(1,1,1)$ vemos que no cumple la condición a pesar de que los vectores sí lo hagan.
		\item $\db{(x,y,z)\in\R^3:x+y+z=0}$
		
		$\begin{array}{l}
			(x_1,y_1,z_1)\in w\\
			(x_2,y_2,z_2)\in W
		\end{array}\longrightarrow(x_1+x_2,y_1+y_2,z_1+z_2)\in W x_1+x_2+y_1+y_2+z_1+z_2=0$
		
		$\begin{aligned}
			\alpha\in\R,(x_1,y_1,z_1)\longrightarrow &\underbrace{\alpha(x_1,y_1,z_1)}\in W\\
			&(\alpha x_1,\alpha y_1,\alpha z_1)\\
		&	\alpha x_1+\alpha y_1+\alpha z_1=\alpha\underbrace{(x_1,y_1,z_1)}_0=0
		\end{aligned}$
	\end{enumerate}
	\item \lb{Comprueba que, en $\R^3$, el subespacio generado por los vectores $(1,2,1)$ y $(6,1,-16)$ coincide con el subespacio generado por los vectores $(-2,7,20)$ y $(4,9,6)$.}
	\item \lb{Usando determinantes, halla una base del subespacio $W$ del ejercicio anterior que esté contenida en el conjunto generador dado.}
	\item \lb{Sean $v_1,v_2,v_3\in\mathbb{K}^n$  con $n\ge3$. Pon un ejemplo de una matriz $A$ tal que los sistemas de ecuaciones $Ax=v_1$ y $Ax=v_2$ tengan solución pero el sistema $Ax=v_3$ no la tenga.}
	
	$\begin{array}{l}
		A\cdot x=\underset{\text{columnas}}{\begin{bmatrix}
			a_1 & \cdots & a_n
	\end{bmatrix}}\begin{bmatrix}
	x_1\\
	\vdots\\
	x_n
\end{bmatrix}=a_1x_1+a_2x_2+\cdots+a_nx_n\\
\underbrace{\begin{bmatrix}
		v_1 & \cdots & 
\end{bmatrix}}_A\underbrace{\begin{bmatrix}
1\\
0\\
\vdots\\
0
\end{bmatrix}}_x=v_1\qquad\underbrace{\begin{bmatrix}
v_1 & v_2 & \cdots
\end{bmatrix}}_{A}\underbrace{\begin{bmatrix}
0 \\
1\\
0\\
\vdots\\
0
\end{bmatrix}}_{x}=v_2\\
A=\underbrace{\begin{bmatrix}
	v_1 & v_2 & 0
\end{bmatrix}}_{A}\underbrace{\begin{bmatrix}
x_1\\
\vdots\\
x_n
\end{bmatrix}}_{x}=v_1x_1+v_2x_2=v_3
	\end{array}$
	\item \lb{Completar la frase: "el vector $b$ pertenece a $\mathrm{Col}(A)$ cuando $\dots\dots$ tiene solución".}
	
	El sistema $Ax=b$ tiene solución
	
	$\begin{array}{l}
		A=\begin{bmatrix}
			a_1 & \cdots & a_n
		\end{bmatrix}\begin{bmatrix}
		x_1\\
		\vdots\\
		x_n
		\end{bmatrix}=a_1x_1+\cdots+a_nx_n=b\\
		\mathrm{Col}(A)=<a_1,a_2,\dots,a_n>
	\end{array}$
	\item \lb{Cierto o falso: "si el vector cero pertenece a $\mathrm{Fil}(A)$, entonces las filas de $A$ son linealmente independientes".}
	
	Si $0\in\mathrm{Fil}(A)\longrightarrow$ las filas de $A$ son linealmente independientes
	
	$\begin{array}{ll}
		A=\begin{bmatrix}
			1 & 0\\
			0 & 1
		\end{bmatrix} & \mathrm{Fil}(A)=<\underbrace{(1,0), (0,1)}_{\text{lin. indep.}}>=\R^2\\
		\begin{tikzpicture}[baseline=(current bounding box.center)]
			\draw (-0.5,0) -- (1.2,0);
			\draw (0,-0.5) -- (0,1.2);
			\draw[lightblue, -latex, line width=1.5pt] (0,0) -- (1,0) node[below] {(1,0)};
			\draw[lightblue, -latex, line width=1.5pt] (0,0) -- (0,1) node[right] {(0,1)};
		\end{tikzpicture} & (0,0)\in\mathrm{Fil}(A)=\R^2
	\end{array}$
	\item \lb{Halla una base del subespacio vectorial $W$ de $\R^5$ generado por los vectores \[ \begin{array}{ccccc}
			(-1,0,0,1,-2) & (2,1,0,-1,2) & (1,3,1,0,-1) & (0,2,1,0,-1) & (3,1,0,-2,4)\\
		\end{array} \]}
	\item \lb{Amplía el conjunto $\{(1,-1,1)\}$ a una base ortonormal de $\R^3$.}
	\item \lb{En $\R^5$, si las ecuaciones del subespacio $U$ son $\left\{\begin{array}{c}
			x+y+z+t+u=0\\
			x-y+z-t+u=0
		\end{array}\right\}$, encuentra las de $U^\perp$.}
	
	$V\equiv\begin{cases}
		x+y+z+t+u=0\\
		x-y+z-y+u=0
	\end{cases}$ ecuaciones \lb{implícitas} $\qquad u^\perp$
	
	Calcular una base de $u$
	
	$\begin{array}{l}
		A=\begin{bmatrix}
			1 & 1 & 1 & 1 & 1 \\
			1 & -1 & 1 & -1 & 1
		\end{bmatrix}\quad \mathrm{rg}(A)=2<5=\text{nº de incógnitas $\longrightarrow$ 3 parámetros }SCI\\
		\begin{rcases}
			u=\alpha\\
			t=\beta\\
			z=\varphi
		\end{rcases}\begin{array}{l}
		x+y=-\varphi-\beta-\alpha\\
		x-y=-\varphi+\beta-\alpha
		\end{array}\longrightarrow\begin{array}{l}
		2x=-2\gamma-2\alpha\longrightarrow x=-\varphi-\alpha\\
		y=-x-\varphi-\beta-\alpha=\cancel{\varphi}+\cancel{\alpha}-\cancel{\varphi}-\beta-\cancel{\alpha}=-\beta
		\end{array}\\
		\begin{cases}
			x=-\alpha-\gamma\\
			y=-\beta\\
			z=\gamma\\
			t=\beta\\
			u=\alpha
		\end{cases}\quad\begin{bmatrix}
		x\\
		y\\
		z \\
		t\\
		u
		\end{bmatrix}=\underbrace{\alpha\begin{bmatrix}
			-1\\
			0\\
			0\\
			0\\
			1
			\end{bmatrix}+\beta\begin{bmatrix}
			0 \\
			-1 \\
			0 \\
			1 \\
			0
		\end{bmatrix}+\gamma\begin{bmatrix}
		-1 \\
		0 \\
		1 \\
		0 \\
		0
	\end{bmatrix}}_{\text{Base de }u}\\
v\in u^\perp\longleftrightarrow v\cdot u=0\:\forall\, u\in U\longleftrightarrow v\cdot u_j=0,\:\{u_1,\dots,u_m\} \text{ base de }U\\
\begin{rcases}
	(x,y,z,t,u)\cdot(-1,0,0,0,1)=0\\
	(x,y,z,t,u)\cdot(0,-1,0,1,0)=0\\
	(x,y,z,t,u)\cdot(-1,0,1,0,0)=0\\
\end{rcases}
	\end{array}$
	
	$\begin{rcases}
		-x+u=0\\
		-y+t=0\\
		-x+z=0
	\end{rcases}$ ecuaciones implícitas de $u^\perp$. 2 parámetros
	
	$\begin{array}{l}
		\underbrace{u}_3\oplus\underbrace{u^\perp}_2=\R^5\qquad\begin{rcases}
		u=\alpha\\
		t=\beta
	\end{rcases}\begin{array}{l}
	x=u=\alpha\\
	y=t=\beta\\
	z=-x=-\alpha
	\end{array}\\
	\begin{bmatrix}
		x \\
		y \\
		z \\
		t \\
		u
	\end{bmatrix}=\underbrace{\alpha\begin{bmatrix}
		1 \\
		0 \\
		1 \\
		0 \\
		1
		\end{bmatrix}+\beta\begin{bmatrix}
		0 \\
		1 \\
		0 \\
		1 \\
		0
	\end{bmatrix}}_{\text{Base de }u^\perp}
	\end{array}$
	
	\item \lb{Halla una base de siguientes subespacios vectoriales de $\R^4$:}
	\begin{enumerate}[label=\color{red}\alph*)]
		\item $\db{\{(x,y,z)|x=y=z=t\}}$
		\item $\db{\{(x,y,z,t)|x+y+z+t=0\}}$
	\end{enumerate} 
	\item \lb{Dados los subespacios \[ \begin{array}{cc}
			U=\{(a,b,a-b,a-2b)\} & W=\{(x,x-y,-x+y,y):x,y\in\R\}
		\end{array} \] calcula la base de cada uno de ellos y determina si están en suma directa.}
	\item \lb{Dada la matriz $U=\begin{bmatrix}
			1 & 0 & 1 & 0 & 1 \\
			0 & 1 & 0 & 1 & 0
		\end{bmatrix}$ calcula:}
	\begin{enumerate}[label=\color{red}\alph*)]
		\item \db{Tres bases distintas de $\mathrm{Col}(U)$.}
		\item \db{Dos bases distintas de $\mathrm{Fil}(U)$.}
	\end{enumerate}
	
	$\begin{array}{l}
		A=\begin{bmatrix}
		1 & 0 & 1 & 0 & 1 \\
		0 & 1 & 0 & 1 & 0
	\end{bmatrix}_{2\times 5}\qquad\mathrm{Fil}(A),\mathrm{Col}(A),\mathrm{nuc}(A),\mathrm{nuc}(A^\intercal )\\
	\mathrm{rg}(A)=\mathrm{dimFil}(A)=\mathrm{dimCol}(A)=2\\
	\mathrm{Fil}(A)=<(1,0,1,0,1), (0,1,0,1,0)>\\
	\mathrm{Col}(A)=<(1,0), (0,1)>
	\end{array}$ 
	
	\includegraphics[width=0.5\linewidth]{"Tareas/Hoja 4/grafica 31-10"}
	
	$\begin{array}{l}
		\underbrace{\mathrm{rg}(A)}_{2}+\underbrace{\mathrm{dimNuc}(A)}_3=5\\
		\begin{rcases}
			x+z+u=0\\
			y+t=0
		\end{rcases}\text{ ecuaciones implícitas de }\mathrm{nuc}(A)\\
		\text{Tomamos 3 parámetros}\\
		\begin{rcases}
			u=\alpha\\
			t=\beta\\
			z=\gamma
		\end{rcases}\begin{array}{l}
		x=-z-u=-\alpha-\gamma\\
		y=-t=-\beta
		\end{array}\\
		\begin{bmatrix}
			x\\
			y\\
			z\\
			t\\
			u
		\end{bmatrix}=\underbrace{\alpha\begin{bmatrix}
			-1\\
			0\\
			0\\
			0\\
			1
			\end{bmatrix}+\beta\begin{bmatrix}
			0\\
			-1\\
			0\\
			1\\
			0
		\end{bmatrix}+\gamma\begin{bmatrix}
		-1\\
		0\\
		1\\
		0\\
		0
	\end{bmatrix}}_{\text{Base de }\mathrm{nuc}(A)}
	\end{array}$
	\item \lb{Sea $S$ un conjunto de 6 vectores de $\R^4$. De las opciones entre paréntesis escoge la(s) correcta(s):}
	\begin{enumerate}[label=\color{red}\alph*)]
		\item \db{$S$ (es) (no es) (no necesariamente es) un conjunto generador de $\R^4$.}
		\item \db{$S$ (es) (no es) (puede ser) un conjunto linealmente independiente.}
		\item \db{Un subconjunto de $S$ con 4 vectores (es) ( no es) (puede ser) una base de $\R^4$}
	\end{enumerate}
	\item \lb{Determina los valores que faltan en la siguiente matriz sabiendo que tiene rango 1: \[ \begin{bmatrix}
			7 & - & - \\
			- & 8 & - \\
			- & 12 & 6 \\
			- & - & 2 \\
			21 & 6 & -
		\end{bmatrix} \]}
	\item \lb{Dadas dos matrices $A$ y $B$, prueba las siguientes afirmaciones:}
	\begin{enumerate}[label=\color{red}\alph*)]
		\item \db{$\mathrm{rg}(A+B)\le\mathrm{rg}(A)+\mathrm{rg}(B)$ (usa que el rango es la dimensión de $\mathrm{Col}(A)$ y lo visto en la página 72 de los apuntes).}
		\item \db{$\mathrm{rg}(AB)\le\mathrm{rg}(A),\mathrm{rg}(B)$ (usa los resultados 4.26 y 4.28 de los apuntes).}
	\end{enumerate}
	\item \lb{Sea $A$ una matriz $m\times n$ de rango $k$. Prueba que existe una matriz $B$ de tamaño $m\times k$ y rango $k$ y una matriz $C$ de tamaño $k\times n$ tales que $A=BC$ (esta factorización se llama \textbf{full rank factorization}). Si $k$ es mucho menor que $n$ y $m$, razona si es mejor, en términos de memoria, almacenar $A$ o almacenar $B$ y $C$.}	
	
	$A=\underset{\text{vectores de tamaño }m}{\begin{bmatrix}
		a_1 & a_2 & \cdots & a_n
	\end{bmatrix}}\quad\mathrm{dimCol}(A)=k\longrightarrow$ Sea $B=\underset{\text{vectores de tamaño }m}{\begin{bmatrix}
	b_1 & \cdots & b_k
\end{bmatrix}}$ una base de $\mathrm{Col}(A)$.

Por tanto, todo vector de $\mathrm{Col}(A)$ se escribe como combinación lineal de los vectores $\{b_1,b_2,\dots,b_k\}$. En particular, las columnas de $A$. \[ \begin{array}{l}
	a_j=c_{ij}b_1+c_{2j}b_2+\cdots+c_{kj}b_k\quad C=[c_{ij}]_{k\times n}\\
	A=\begin{bmatrix}
		a_1 & \cdots & a_n
	\end{bmatrix}=\begin{bmatrix}
	b_1 & \cdots & b_k
	\end{bmatrix}_{n\times k}\cdot \left[\begin{array}{cc:c}
	c_{11}&c_{12}  & c_{1n} \\
	c_{21}& c_{22} & c_{2n} \\
	\vdots& \vdots & \vdots \\
	c_{k1}& c_{k2} &c_{kn} 
	\end{array}\right]_{k\times n}
\end{array}\]
\Ej

$\begin{array}{l}
	m=10^2 \\
	n=10^2 \\
	k=10 \\
\end{array}\begin{array}{l}
 \text{$A$ necesita almacenar }10^2\cdot10^2=10^4=10,000\text{ datos}\\
\begin{rcases}
	\text{$B$ necesita almacenar $10^2\cdot10=10^3=1,000$ datos}\\
\text{$C$ necesita almacenar $10^2\cdot10=10^3=1,000$ datos}
\end{rcases}2,000\text{ datos}
\end{array}$
	
	\item \lb{Calcula la "full rank factorization" para la matriz \[ \begin{bmatrix}
			1 & 0 & 2 & 3 \\
			-1 & -3 & 1 & 0
		\end{bmatrix} \]}
	
	$\begin{array}{l}
		A=\begin{pmatrix}
			1 & 0 & 2 & 3 \\
			-1 & -3 & 1 & 0
		\end{pmatrix}\qquad\mathrm{rg}(A)=2\\
		B=\underset{\text{base de $\mathrm{Col}(A)$}}{\begin{bmatrix}
				1 & 0\\
				-1 & -3
		\end{bmatrix}_{2\times2}}\qquad C=\begin{bmatrix}
		1 & 0 & 2 & 3 \\
		0 & 1 & -1 & -1
	\end{bmatrix}\\
	\begin{bmatrix}
		1\\
		-1
	\end{bmatrix}=c_{11}\begin{bmatrix}
	1\\
	-1
	\end{bmatrix}+c_{21}\begin{bmatrix}
	0\\
	-3
	\end{bmatrix}\\
	\begin{bmatrix}
		0\\
		-3
	\end{bmatrix}=c_{12}\begin{bmatrix}
	1\\
	-1
	\end{bmatrix}+c_{22}\begin{bmatrix}
	0\\
	-3
	\end{bmatrix}\\
	\begin{bmatrix}
		2\\
		1
	\end{bmatrix}=c_{13}\begin{bmatrix}
	1\\
	-1
	\end{bmatrix}+c_{23}\begin{bmatrix}
	0\\
	-3
	\end{bmatrix}\\
	\begin{rcases}
		2=c_{13}\\
		1=-c_{13}-3c_{23}
	\end{rcases}\begin{array}{l}
	\\
	3c_{23}=-1-c_{13}=-3\longrightarrow c_{23}=-1
	\end{array}\\
	\begin{bmatrix}
		3\\
		0
	\end{bmatrix}=c_{14}\begin{bmatrix}
	1\\
	-1
	\end{bmatrix}+c_{24}\begin{bmatrix}
	0\\
	-3
	\end{bmatrix}\\
	\begin{rcases}
		3=c_{14}\\
		0=-c_{14}-3c_{24}
	\end{rcases}\begin{array}{l}
	\\
	3c_{24}=-c_{14}=3\longrightarrow c_{24}=-1
	\end{array}
	\end{array}$
	
	\bu{Comprobación}
	\[ \underbrace{\begin{bmatrix}
			1 & 0\\
			-1 & -3
	\end{bmatrix}}_{B}\cdot\underbrace{\begin{bmatrix}
	1 & 0 & 2 & 3 \\
	0 & 1 & -1 & -1
\end{bmatrix}}_C=\underbrace{\begin{bmatrix}
1 & 0 & 2 & 3 \\
-1 & -3 & 1 & 0
\end{bmatrix}}_A \]
	\item \lb{Utiliza la "full rank factorization" y el ejercicio 15 para ver que una matriz $A$ tiene rango $k$ si y sólo si se puede expresar como suma de $k$ matrices de rango 1, pero no se puede expresar como suma de menos de $k$ matrices de rango 1. Expresa la matriz del ejercicio anterior como suma de 2 matrices de rango 1.}
	\item \lb{Se consideran las matrices \[ \begin{array}{cc}
			A=\begin{bmatrix}
				1 & 3 & 2 \\
				0 & 1 & 1 \\
				1 & 3 & 2
			\end{bmatrix} & B=\begin{array}{ccc}
			1 & 3 & 2 \\
			0 & 1 & 1 \\
			0 & 0 & 0
			\end{array}
		\end{array} \] en donde $B$ se obtiene restando la fila uno a la fila tres. ¿Qué relación hay entre los cuatro subespacios fundamentales de las dos matrices? Calcula una base para cada uno de ellos.}
		\item \lb{Pon un ejemplo de una matriz cuadrada $A$ en la que $\mathrm{Col}(A)=\mathrm{Fil}(A)$ y otro en el que $\mathrm{Col}(A)\neq\mathrm{Fil}(A)$}
        \item \lb{Sea $A$ una matriz $10\times 10$ que cumple $A^2=0$. Prueba que $\mathrm{Col}(A)\subseteq\mathrm{Nuc}(A)$ y, en consecuencia, el rango de $A$ es menor o igual a 5.}
        
        Sea $v\in\mathrm{Col}(A)\xRightarrow{?} Av=0$\\
        Como $v\in\mathrm{Col} (A)\longrightarrow\exists\,x\in\R^{10}$ tal que \[ v=Ax=\begin{bmatrix}
        	a_1& \cdots & a_{10}
        \end{bmatrix} \cdot\begin{bmatrix}
        x_1\\
        \vdots\\
        x_{10}
        \end{bmatrix}=a_1x_1+\cdots+a_{10}x_{10}\]
        $Av=A\cdot(Ax)=A^2x=0\longrightarrow v\in\mathrm{nuc}(A)$
        
        \begin{tikzpicture}
        	\node[red,draw=red, fill=red!10, rectangle, line width=1.5]  {\underline{Nota:} $\underset{\le 5}{\mathrm{rg}(A)}+\mathrm{dimnuc}(A)=10,\:\mathrm{rg}(A)=\mathrm{dimCol}(A)$};
        \end{tikzpicture}
        
        \item \lb{Sea $A$ una matriz cuadrada invertible. Calcula una base para cada uno de los subespacios fundamentales de las matrices $A$ y $B=\begin{bmatrix}
        		A & A
        	\end{bmatrix}$.}
        \item \lb{Sin calcular la matriz $A$, encuentra bases para sus cuatro subespacios fundamentales \[ A=\begin{bmatrix}
        		1 & 0 & 0 \\
        		6 & 1 & 0 \\
        		9 & 8 & 1
        	\end{bmatrix}\:\begin{bmatrix}
        	1 & 2 & 3 & 4 \\
        	0 & 1 & 2 & 3 \\
        	0 & 0 & 1 & 2
        	\end{bmatrix} \]}
        \item \lb{Halla una base del subespacio $\R^4$ dado por el núcleo de $A$, donde \[ A=\begin{bmatrix}
        		3 & 0 & 1 & 2 \\
        		1 & 1 & 0 & 1
        	\end{bmatrix} \]Comprueba que el vector $(1,1,1-2)$ pertenece a $\mathrm{Nuc}(A)$ y calcula sus coordenadas respecto de la base obtenida.}
        \[ \begin{bmatrix}
        	3 & 0 & 1 & 2 \\
        	1 & 1 & 0 & 1
        \end{bmatrix}\cdot\begin{bmatrix}
        x\\
        y\\
        z\\
        t
        \end{bmatrix}=\begin{bmatrix}
        3x+z+2t\\
        x+y+t
        \end{bmatrix}=\begin{bmatrix}
        0\\
        0
        \end{bmatrix}\]
        $\begin{rcases}
        	2x+z+2t=0\\
        	x+y+t=0
        \end{rcases}\qquad\begin{array}{l}
        \mathrm{rg}(A)=2\\
        \mathrm{rg}(A)+\mathrm{dimNuc}(A)=4\longrightarrow\mathrm{dimNuc}(A)=2
        \end{array}$
        
        Tomamos 2 parámetros:
        
        $\begin{array}{l}
        	 \begin{rcases}
        	t=\alpha\\
        	z=\beta
        \end{rcases}\longrightarrow\begin{array}{l}
        3x=-2t-z=-2\alpha-\beta\longrightarrow x=-\dfrac{2}{3}\alpha-\dfrac{1}{3}\beta\\
        y=-t-x=-\alpha+\dfrac{2}{3}\beta=-\dfrac{1}{3}\alpha+\dfrac{1}{3}\beta
        \end{array}\\
        \begin{bmatrix}
        	x\\
        	y\\
        	z\\
        	t
        \end{bmatrix}=\alpha\underset{\downarrow\times3}{\begin{bmatrix}
        -\frac{2}{3}\\
        -\frac{1}{3}\\
        0\\
        1
        \end{bmatrix}}+\beta\underset{\downarrow\times3}{\begin{bmatrix}
        -\frac{1}{3}\\
        \frac{1}{3}\\
        1\\
        0
        \end{bmatrix}}\\
        B_{\mathrm{nuc}(A)}=\{(-2,-1,0,3),(-1,1,3,0)\}\\
        \begin{array}{cccc}
        	3 & 0 & 1 & 2 \\
        	1 & 1 & 0 & 1
        \end{array}\cdot\begin{bmatrix}
        1\\
        1\\
        1\\
        -2
        \end{bmatrix}=\begin{bmatrix}
        3+1-4\\
        1\cdot1+1\cdot1-2
        \end{bmatrix}=\begin{bmatrix}
        0\\
        0
        \end{bmatrix}\longrightarrow(1,1,1,1-2)\in\mathrm{nuc}(A)\\
        (1,1,1,-2)=\alpha(-2,-1,0,3)+\beta(-1,1,3,0)\\
        \begin{rcases}
        	1=-2\alpha-\beta\\
        	1=\alpha+\beta\\
        	1=3\beta\\
        	-2=3\alpha
        \end{rcases}\qquad\begin{array}{l}
        \alpha=-\dfrac{2}{3}\\
        \beta=\dfrac{1}{3}
        \end{array}\\
        \left[1,1,1,-2\right]_{B_{\mathrm{nuc}(A)}}=\left(-\dfrac{2}{3},\dfrac{1}{3}\right)
        \end{array}$
        \item \lb{Dadas las siguientes bases de $\R^4$ \[ \begin{array}{c}
        		\mathcal{B}_1=\{v_1=(1,1,0,0),v_2=(0,1,0,0),v_3=(0,0,1,1), v_4=(0,0,0,1)\}\\
        		\mathcal{B}_2=\{w_1=(1,2,0,0), w_2=(0,1,2,-1), w_3=(0,1,1,1), w_4=(0,1,2,0)\}
        	\end{array} \] encuentra la matriz de cambio de base de $\mathcal{B}_1$ a $\mathcal{B}_2$ y la de l cambio inverso. ¿Qué coordenadas tiene el vector $3w_1-w_3+2w_2$ con respecto a la base $\mathcal{B}_1$?}
        
        $M_{B_1\to B_2}=\left[\begin{array}{c|c}
        	1 & \\
        	-2 & \\
        	-2 & \\
        	3 & 
        \end{array}\right]\qquad M_{B_2\to B_1}=\begin{bmatrix}
        M_{B_1\to B_2}
        \end{bmatrix}^{-1}$
        
        $\begin{aligned}
        	v_1&=a_{11}w_1+a_{21}w_2+a_{31}w_3+a_{41}w_4\\
        	&=a_{11}(1,2,0,0) +a_{21}(0,1,2,-1)+_{31}(0,1,1,1)+a_{41}(0,1,2,0)\\
        	&=(1,1,0,0)\\
        	&\begin{rcases}
        		1=a_{11}\cdot1\\
        		1=2a_{11}+a_{21}+a_{31}+a_{41}\\
        		0=2a_{21}+a_{31}+2a_{41}\\
        		0=-a_{21}+a_{31}
        	\end{rcases}\begin{rcases}
        	-1=\overset{-2}{a_{21}}+\overset{-2}{a_{31}}+a_{41}\\
        	0=2a_{21}+a_{31}+2a_{41}\\
        	0=-a_{21}+a_{31}
        	\end{rcases}\xRightarrow{F_1\to2F_1}\begin{array}{rcl}
        	-2&=&2a_{21}+2a_{31}+2a_{41}\\
        	0&=&2a_{21}+a_{31}+2a_{41}\\ \hline
        	-2&=&a_{31}
        	\end{array}
        \end{aligned}$
        
        $\begin{array}{l}
        	3v_1-v_3+v_2=\begin{bmatrix}
        		3 & 1 & -1 & 0
        	\end{bmatrix}_{B_1}\\
        	\left[3v_1-v_3+v_2\right]_{B_2}=M_{B_1\to B_2}[3v_1-v_3+v_2]_{B_1}=\left[\begin{array}{c|c}
        		1 & \\
        		-2 & \\
        		-2 & \\
        		3 & 
        	\end{array}\right]\cdot\begin{bmatrix}
        	3\\
        	1\\
        	-1\\
        	0
        	\end{bmatrix}
        \end{array}$
        \item \lb{Calcula una base y unas ecuaciones implícitas de $U+V$, donde \[ \begin{array}{c}
        		U=\left\langle(1,2,3,4,5,), (5,4,3,2,1),(-1,0,1,2,3)\right\rangle \\
        		V=\left\{(x,y,z,r,s):\begin{array}{l}
        			2x+r=2y\\
        			3x+s=2y
        		\end{array}\right\}
        	\end{array} \]}
        
        $\begin{array}{l}
        	V\equiv\begin{cases}
        	2x+r=2y\\
        	3x+s=2y
        \end{cases}\longrightarrow\begin{rcases}
        2x-2y+r=0\\
        3x-2y+s=0
        \end{rcases}\\
        \underbrace{\begin{array}{ccccc}
        		3 & -2 & 0 & 1 & 0 \\
        		3 & -2 & 0 & 0 & 1
        \end{array}}_A\cdot\begin{bmatrix}
        x\\
        y\\
        z\\
        t
    \end{bmatrix}=\begin{bmatrix}
    0\\
    0
    \end{bmatrix}\\
    \mathrm{rg}(A)=2<3=\text{ nº de incógnitas }\longrightarrow S.C.I
        \end{array}$
        
        Tomamos 3 parámetros
        
        $\begin{array}{l}
        	\begin{rcases}
        		s=\alpha\\
        		r=\beta\\
        		z=\gamma
        	\end{rcases}\qquad\begin{rcases}
        	2x-2y=-\beta\\
        	3x-2y=\alpha
        	\end{rcases}\qquad\begin{array}{l}
        	x=\alpha+\beta\\
        	y=\dfrac{1}{2}(3\alpha+3\beta-\alpha)=\alpha+\dfrac{3}{2}\beta
        	\end{array}\\
        	\begin{bmatrix}
        		x\\
        		y\\
        		z\\
        		r\\
        		s
        	\end{bmatrix}=\alpha\begin{bmatrix}
        	1\\
        	1\\
        	0\\
        	0\\
        	1
        	\end{bmatrix}+\beta\begin{bmatrix}
        	1\\
        	\frac{3}{2}\\
        	0\\
        	1\\
        	0
        	\end{bmatrix}+r\begin{bmatrix}
        	0\\
        	0\\
        	1\\
        	0\\
        	0
        	\end{bmatrix}\\
        	V=<(1,1,0,0,1),\left(1,\dfrac{3}{2},0,1,0\right),(0,0,1,0,0)>\\
        	U+V=<(1,2,3,4,5),(5,4,3,2,1),(-1,0,1,2,3),(1,1,0,0,1),\left(1,\dfrac{3}{2},0,1,0\right),(0,0,1,0,0)>
        \end{array}$
        
        Base de $U+V$ se obtiene haciendo operaciones elementales sobre la matriz \[ A=\begin{bmatrix}
        	1 & 5 & -1 & 1 & 1 & 0 \\
        	2 & 4 & 0 & 1 & \frac{3}{2} & 0 \\
        	3 & 3 & 1 & 0 & 0 & 1 \\
        	4 & 2 & 2 & 0 & 1 & 0 \\
        	5 & 1 & 3 & 1 & 0 & 0
        \end{bmatrix} \] de igaul forma a como lo hacemos para calcular el rango. Las columnas que resulten con pivotes no nulos constituyen una base de $U+V$.
        
        Las ecuaciones implícitas los podemos obtener a partir de las paramétricas, eliminando los parámetros. También se puede hacer con determinantes.
\end{enumerate}
